%==============================================================================
\chapter{التصميم (نماذج UML)}
\label{chap:conception}
%==============================================================================
يتضمّن ملف التصميم مخطّطات \textenglish{UML} التالية. كلٌّ منها موصوف بمحتواه
المنتظر لتوجيه النمذجة.
\begin{xltabular}{\textwidth}{>{\bfseries}p{3.8cm} X}
\toprule
\textbf{المخطّط} & \textbf{المحتوى المنتظر} \\
\midrule
\endhead
حالات الاستعمال & الفاعلون (مربّي النحل، عون، مشرف، مسؤول، مدير نظام، تطبيق خارجي) والحالات الرئيسية (انظر الشكل \ref{fig:uc}). \\
الأصناف & كيانات قاموس المعطيات وعلاقاتها (مربّي النحل 1 إلى n مزرعة، مزرعة 1 إلى n موقع، موقع 1 إلى n خلية، خلية 1 إلى n جسم أو عاسلة، جسم أو عاسلة 1 إلى n إطار) وطرق الخدمة. \\
الكائنات & لقطة لموقع يحتوي على ثلاث خلايا مُعمَّرة بإطاراتها. \\
التسلسل & سير إنجاز زيارة (العون، واجهة JSP، Servlet، EJB، JDBC، قاعدة المعطيات). \\
النشاط & سيرورة تخطيط زيارة والمصادقة عليها وتنفيذها وإغلاقها. \\
آلة الحالات & دورة حياة خلية (مُنشأة، مُعمَّرة، نشطة، في تقسيم، في جمع، مُغلقة). \\
التفاعل أو التعاون & السيناريو نفسه للتسلسل في شكل تعاون بين الكائنات. \\
الحزمة & طبقات العرض والتحكم والأعمال والوصول إلى المعطيات والتدويل والإعداد وخدمات الويب. \\
المكوّن & المكوّنات التطبيقية (تطبيق ويب، خدمة OAuth، خدمة API، وحدة i18n، وحدة الإعداد، قاعدة المعطيات). \\
النشر & التوزيع على المتصفّح وخادم تطبيقات J2EE ونظام إدارة قواعد المعطيات وخدمة Google OAuth والمستشعرات. \\
\bottomrule
\end{xltabular}

\section{مخطّط الأصناف}
يصف مخطّط الأصناف، المعروض بمحاذاة أفقية، الكيانات المهنية وخصائصها وطريقة الخدمة
\texttt{getZummHoneyActualQuantity} والعلاقات بتعدّدياتها.
\begin{figure}[htbp]
\centering
\includegraphics[width=\linewidth]{diagramme_classes}
\caption{مخطّط أصناف نظام Zümm (محاذاة أفقية).}
\label{fig:classes}
\end{figure}

\section{العرض الطبقي (الحزمة)}
\begin{figure}[htbp]
\centering
\begin{tikzpicture}[
  couche/.style={draw=mielfonce,fill=grisdoux,thick,minimum width=10cm,minimum height=8mm,font=\small}
]
\node[couche] (p) at (0,4) {العرض: JSP وJavaScript};
\node[couche] (c) at (0,3) {التحكم: Servlets (MVC)};
\node[couche] (m) at (0,2) {الأعمال: EJB};
\node[couche] (d) at (0,1) {الوصول إلى المعطيات: JDBC};
\node[couche] (b) at (0,0) {قاعدة معطيات علائقية};
\node[draw=mielfonce,thick,fill=miel!15,minimum height=8mm,font=\footnotesize,right=3mm of c,xshift=4.7cm] (s) {i18n XML / ConfigZumm.ini / خدمات الويب};
\draw[-{Latex[length=2mm]},mielfonce] (p)--(c);
\draw[-{Latex[length=2mm]},mielfonce] (c)--(m);
\draw[-{Latex[length=2mm]},mielfonce] (m)--(d);
\draw[-{Latex[length=2mm]},mielfonce] (d)--(b);
\end{tikzpicture}
\caption{بنية متعددة الطبقات في طبقات ضعيفة الاقتران.}
\label{fig:couches}
\end{figure}

\section{عرض النشر}
\begin{figure}[htbp]
\centering
\begin{tikzpicture}[
  noeud/.style={draw=mielfonce,fill=grisdoux,thick,minimum width=3cm,minimum height=1.1cm,font=\footnotesize,align=center},
  fl/.style={-{Latex[length=2mm]},mielfonce,thick}
]
\node[noeud] (nav) {المتصفّح\\(عميل ويب)};
\node[noeud,right=1.6cm of nav] (srv) {خادم J2EE\\(Servlets، EJB)};
\node[noeud,right=1.6cm of srv] (bd) {نظام إدارة\\قواعد معطيات};
\node[noeud,below=1cm of srv] (oauth) {Google OAuth};
\node[noeud,below=1cm of bd] (capt) {مستشعرات\\(الشقّ الآلي)};
\draw[fl] (nav)--(srv) node[midway,above,font=\scriptsize]{HTTPS};
\draw[fl] (srv)--(bd) node[midway,above,font=\scriptsize]{JDBC};
\draw[fl] (srv)--(oauth) node[midway,right,font=\scriptsize]{SSL};
\draw[fl] (capt)--(srv);
\end{tikzpicture}
\caption{مخطّط نشر النظام.}
\label{fig:deploiement}
\end{figure}

\section{مخطّط التسلسل (UC2)}
يفصّل مخطّط التسلسل أدناه حالة الاستعمال UC2، إنجاز زيارة وملء التقرير. يُظهر
المرور عبر الطبقات (العون، الواجهة، servlet، EJB، قاعدة المعطيات) إضافةً إلى
الحالتين المتّصلة وغير المتّصلة مع مزامنة مؤجّلة.
\begin{figure}[htbp]
\centering
\includegraphics[width=0.92\linewidth]{diagramme_sequence_uc2}
\caption{مخطّط تسلسل حالة الاستعمال UC2 (إنجاز زيارة).}
\label{fig:sequence-uc2}
\end{figure}

\section{تسلسل المصادقة (Google OAuth)}
يتّبع مسار المصادقة تدفّق \emph{Authorization Code} لـ OAuth~2.0. لا ينقل
المتصفّح أبدًا السرّ الخاص بالعميل: يتمّ تبادل رمز التفويض بالرمز المميّز للوصول
\textbf{من خادم إلى خادم}، عبر قناة SSL مُوثَّقة بشهادة X.509. تُفتَح جلسة
التطبيق بعد ذلك بواسطة كعكة آمنة (\texttt{HttpOnly}، \texttt{Secure})، ويُربَط
الملف الذي تُرجعه Google بعون \texttt{Agent} موجود (أو يُنشئ واحدًا في انتظار
التأهيل). في العرض، يُتوقَّع تراجع إلى مصادقة محلية عند عدم توفّر المزوّد.
\begin{figure}[htbp]
\centering
\begin{tikzpicture}[font=\scriptsize,
  msg/.style={-{Latex[length=1.6mm]},mielfonce,thick},
  rmsg/.style={-{Latex[length=1.6mm]},mielfonce,thick,dashed}
]
\foreach \x/\name in {0/{العون (المتصفّح)}, 6/{خادم Zümm}, 12/{Google OAuth}}{
  \node[draw=mielfonce,fill=grisdoux,thick,align=center,minimum height=7mm,inner xsep=6pt] at (\x,0) {\name};
  \draw[mielfonce,dashed] (\x,-0.45) -- (\x,-9.9);
}
\draw[msg]  (0,-1.0) -- (6,-1.0)  node[midway,above]{1. الوصول إلى صفحة محمية};
\draw[rmsg] (6,-1.7) -- (0,-1.7)  node[midway,above]{2. 302: إعادة توجيه (client\_id، scope، state)};
\draw[msg]  (0,-2.4) -- (12,-2.4) node[midway,above]{3. المصادقة + الموافقة};
\draw[rmsg] (12,-3.1) -- (0,-3.1) node[midway,above]{4. 302: إعادة توجيه + code + state};
\draw[msg]  (0,-3.8) -- (6,-3.8)  node[midway,above]{5. GET /callback (code، state)};
\draw[msg]  (6,-4.5) -- (12,-4.5) node[midway,above]{6. POST /token (code، secret) --- SSL / X.509};
\draw[rmsg] (12,-5.2) -- (6,-5.2) node[midway,above]{7. access\_token + id\_token (JWT)};
\draw[msg]  (6,-5.9) -- (12,-5.9) node[midway,above]{8. GET /userinfo (access\_token)};
\draw[rmsg] (12,-6.6) -- (6,-6.6) node[midway,above]{9. الملف (email، الاسم)};
\node[draw=mielfonce,fill=miel!12,thick,align=center,inner sep=3pt] at (6,-7.5)
  {10. التحقّق من \texttt{state}، ربط العون،\\فتح الجلسة (RBAC)};
\draw[rmsg] (6,-8.6) -- (0,-8.6) node[midway,above]{11. 302: التطبيق + كعكة جلسة آمنة};
\end{tikzpicture}
\caption{تسلسل مصادقة OAuth~2.0 (Authorization Code) مع Google.}
\label{fig:seq-oauth}
\end{figure}

\section{مخطّط الكائنات}
يعرض مخطّط الكائنات لقطة لموقع يستضيف ثلاث خلايا، مع تفصيل تركيب خلية إلى جسم
وعاسلة وإطارات.
\begin{figure}[htbp]
\centering
\includegraphics[width=0.7\linewidth]{diagramme_objets}
\caption{مخطّط الكائنات، لقطة لموقع تربية نحل.}
\label{fig:objets}
\end{figure}

\section{مخطّط النشاط}
يُنمذج مخطّط النشاط سيرورة تخطيط زيارة والمصادقة عليها وإنجازها، مع البديلين
المتّصل وغير المتّصل.
\begin{figure}[htbp]
\centering
\includegraphics[width=0.7\linewidth]{diagramme_activite}
\caption{مخطّط نشاط سيرورة الزيارة.}
\label{fig:activite}
\end{figure}

\section{مخطّط آلة الحالات}
يصف مخطّط آلة الحالات دورة حياة خلية، من إنشائها إلى إغلاقها.
\begin{figure}[htbp]
\centering
\includegraphics[width=0.85\linewidth]{diagramme_etat}
\caption{آلة حالات دورة حياة خلية.}
\label{fig:etat}
\end{figure}

\section{مخطّط التعاون}
يستعيد مخطّط التعاون سيناريو UC2 في شكل رسائل مرقّمة مُتبادَلة بين الكائنات.
\begin{figure}[htbp]
\centering
\includegraphics[width=\linewidth]{diagramme_collaboration}
\caption{مخطّط تعاون حالة الاستعمال UC2.}
\label{fig:collaboration}
\end{figure}

\section{مخطّط المكوّنات}
يُظهر مخطّط المكوّنات تنظيم المكوّنات التطبيقية وواجهاتها وتبعياتها، ومنها واجهة
الخدمة التي تعرض كمية العسل.
\begin{figure}[htbp]
\centering
\includegraphics[width=\linewidth]{diagramme_composant}
\caption{مخطّط مكوّنات نظام Zümm.}
\label{fig:composant}
\end{figure}
